%=============================================================================!
% 15.3  BLOCK AVERAGES                                                        !
%=============================================================================!
\section{Block averages}
\label{sec:cv-blocks}

The estimate of $g$ rests on a rule for where the autocorrelation function
ends, and a slow tail that the rule cuts off would leave the error bar too
small. Block averaging finds the error of the mean without estimating the
autocorrelation function at all, and shows on a plot whether the run was
long enough to find it\cite{Flyvbjerg1989}.

\subsection*{Blocks of samples}

We cut the $n$ samples of a run into $n_{\mathrm b}$ consecutive blocks of
$b$ samples each, $n = n_{\mathrm b}b$ (a remainder at the end is dropped),
and average each block, giving the block means $\bar A_1, \dots, \bar
A_{n_{\mathrm b}}$, whose own mean is $\bar A$. Each block mean is the mean of
$b$ correlated samples, so by \cref{eq:cv-var} with $b$ in place of $n$ it
has the variance $g_b\sigma_A^2/b$, with
\begin{equation*}
   g_b = 1 + 2\sum_{k=1}^{b-1}\Bigl(1 - \frac kb\Bigr)
   \operatorname{corr}(k) .
\end{equation*}
Two neighbouring blocks share only the correlation between the samples
near their common edge, so blocks much longer than the run's memory have
nearly independent means. The mean of $n_{\mathrm b}$ independent block
means then has the variance of one divided by $n_{\mathrm b}$
(\cref{eq:cv-mean}), estimated from their own spread: the error of $\bar
A$ is $s_{\mathrm b}/\sqrt{n_{\mathrm b}}$, with $s_{\mathrm b}^2$ the sum of
the squared deviations of the block means from $\bar A$ divided by
$n_{\mathrm b} - 1$. No autocorrelation function enters.

How long the blocks must be follows from $g_b$. Splitting the sum,
\begin{equation*}
   g_b = 1 + 2\sum_{k=1}^{b-1}\operatorname{corr}(k)
   - \frac2b\sum_{k=1}^{b-1}k\operatorname{corr}(k)
   \approx g - \frac2b\sum_{k=1}^{\infty}k\operatorname{corr}(k) ,
\end{equation*}
since for blocks much longer than the memory both sums have already
reached their limits. For the exponential $\operatorname{corr}(k) = c^k$
of \cref{eq:cv-ou}, the second sum is $c/(1 - c)^2$ (\cref{ex:cv-ksum}),
which is $(g^2 - 1)/4$, since $g^2 - 1 = [(1 + c)^2 - (1 - c)^2]/(1 - c)^2
= 4c/(1 - c)^2$; for $g$ well above 1 it is close to $g^2/4$, so $g_b
\approx g - g^2/(2b)$. The blocks estimate the error
$\sqrt{g_b\sigma_A^2/(n_{\mathrm b}b)} = \sqrt{g_b\sigma_A^2/n}$, against
the true $\sqrt{g\sigma_A^2/n}$; the part missing is the correlation
between neighbouring blocks. Their ratio is $\sqrt{1 - x}$ with $x =
g/(2b)$, and a small relative change changes a square root by half as
much, so blocks of $b$ samples give an error too small by the fraction
\begin{equation}
   1 - \sqrt{g_b/g} \approx \frac{g}{4b} .
   \label{eq:cv-short}
\end{equation}
For the series with $g = 39$ the exact shortfall is 0.102 for $b = 100$
and 0.051 for $b = 195$, against 0.097 and 0.050 from \cref{eq:cv-short}.
Blocks of $5g$ samples miss about 5\% of the error when the memory is
exponential, a factor that the rest of the chapter uses. A memory that
falls quickly at first and then fades slowly, as the liquid's does
(\cref{fig:cv-acf}), has a larger $\sum_kk\operatorname{corr}(k)$ for the
same $g$, 1.7 to 2.0 times $(g/2)^2$ for $U$, $K$ and $P$, and blocks of
$5g$ samples then miss 0.107, 0.089 and 0.090 of the error, within the
error of an estimate from twenty blocks.

\subsection*{The plateau}

Flyvbjerg and Petersen double the block length at each stage, from 1
sample to as many as leave a few blocks, and plot the error against the
length\cite{Flyvbjerg1989}. Short blocks give the error of independent
samples, too small by $\sqrt g$; as the blocks grow past the memory the
estimate rises and levels off at $\sqrt{g\sigma_A^2/n}$. The level part is
the plateau, and its height is the error of the mean. Its own precision
is that of an estimate of a standard deviation from $n_{\mathrm b}$
samples: for nearly Gaussian block means the relative error of $s_{\mathrm
B}$ is $1/\sqrt{2(n_{\mathrm b} - 1)}$ (\cref{ex:cv-error-error}), 0.24 for
10 blocks and 0.16 for 20. The longest blocks, few in number, scatter.
The function \texttt{block\_average} of \texttt{mdlab.analysis.stats}
returns the error and its own error at each length:

\begin{code}
def block_average(series: ArrayLike, min_blocks: int = 4) -> dict:
    a = np.asarray(series, dtype=float)
    lengths, errors, error_errors, counts = [], [], [], []
    b = 1
    while len(a) // b >= min_blocks:
        n_b = len(a) // b
        means = a[: n_b * b].reshape(n_b, b).mean(axis=1)
        se = means.std(ddof=1) / np.sqrt(n_b)
        lengths.append(b)
        errors.append(se)
        error_errors.append(se / np.sqrt(2 * (n_b - 1)))
        counts.append(n_b)
        b *= 2
    return {"length": np.array(lengths), "error": np.array(errors),
            "error_error": np.array(error_errors),
            "blocks": np.array(counts)}
\end{code}

\noindent \Cref{fig:cv-blocks}(a) applies it to the series of
\cref{sec:cv-correlation}. With $2^{17}$ samples, 3361 independent ones,
the error rises from 0.0028 for single samples to a plateau of 0.0164 to
0.0168 for blocks of 256 to 1024 samples, against the exact 0.0172. With
$2^{10}$ samples, 26 independent ones, it rises to 0.182 for blocks of 64
and falls back to 0.131 for blocks of 256, of which only four remain,
while the exact value is 0.193: there is no plateau, and a run that shows
none is too short for its own error bar to be known. The criterion of
\cref{sec:cv-protocol} asks for at least 100 independent samples, enough
for twenty blocks of $5g$ samples and an error known to about a sixth.

\begin{figure}[htb]
\centering
\includegraphics{ch15_convergence/blocks}
\caption{Block averages. (a) The series of \cref{eq:cv-ou} with $c =
0.95$: the error of the mean from blocks of $b$ samples, for $2^{17}$
samples (teal) and $2^{10}$ (ochre), with the error of each estimate,
against the exact errors (dashed). (b) The liquid under CSVR, \SI{2}{\nano
\second}: the error of the mean of $U$, $K$ and $P$ from blocks, divided by
$\sqrt{g s^2/n}$. (c) The same run cut into twenty runs of
\SI{100}{\pico\second}: the mean of $U$ in each, less that of the whole
run, with its own error bar from that piece alone.}
\label{fig:cv-blocks}
\end{figure}

\subsection*{The liquid, and a test against independent runs}

For the liquid of \cref{fig:cv-acf}, blocks of \SI{20}{\pico\second} give
1.03, 1.01 and 1.00 times the errors found from $g$ for $U$, $K$ and $P$,
and blocks of \SI{41}{\pico\second} 0.98, 0.99 and 0.97
(\cref{fig:cv-blocks}(b)): the rule of \cref{sec:cv-correlation} missed no
slow tail that blocks known to 7\% could show. The longest blocks,
\SI{164}{\pico\second} and only twelve of them, give 1.20 to 1.28, about
one of their own errors, a fifth, above 1.

The sharper test is whether an error bar from one run predicts how far
other runs fall from it. Cutting the \SI{2}{\nano\second} run into twenty
pieces of \SI{100}{\pico\second}, whose means are nearly independent
(successive means have the correlation 0.06), gives twenty runs to set
against one another (\cref{fig:cv-blocks}(c)). Their means of $U$ scatter
by \SI{0.0241}{\electronvolt}, while the error bar that each piece finds
for itself averages \SI{0.0196}{\electronvolt}, from 0.0128 to
\SI{0.0278}{\electronvolt}; for $K$ the figures are 0.0351 and
\SI{0.0312}{\electronvolt}, and for $P$ 0.00201 and
\SI{0.00172}{\giga\pascal}. Eleven of the twenty bars for $U$ hold the mean
of the whole run, where one standard error should hold it in 68\% of
runs, 13.7 of twenty on average, a count that itself spreads by 2.1, since
each piece holds the mean or not, a yes or no of variance $0.68\times0.32$,
and the twenty variances add. A piece of \SI{100}{\pico\second} holds 68
independent samples of $U$, and the $g$ it finds for itself ranges from 64
to 228, on either side of the 146 of the whole run: an error bar from so
short a run is uncertain by about a fifth (the twenty bars for $U$ spread
by 0.20 of their mean), and on average a little low.

\begin{keyidea}
Block means of blocks much longer than the run's memory are nearly
independent, and their spread gives the error of the mean with no
autocorrelation function. Plotted against the block length, the error
rises to a plateau; blocks of $5g$ samples miss about 5\% of it for an
exponential memory, and the
plateau's own relative error is $1/\sqrt{2(n_{\mathrm b} - 1)}$. A run
with no plateau is too short to know its error bar.
\end{keyidea}

\notebook{15}{cut a series into blocks and find the plateau}

\begin{selfcheck}
A run of 20\,000 samples has $g = 100$ and an exponential memory. How long
must its blocks be to miss about 5\% of the error, and how many such
blocks does the run hold?
\end{selfcheck}
